%=====================================================================
%  05-QUANTUM.TEX — PRINSIP DASAR KUANTUM
%=====================================================================
%  Contoh penulisan notasi: \ket{...}, \hat{U}, \rho, bmatrix.
%  Typo pada versi lama (\mathcal{p}_i, "qubt", "partiker", "dnegan")
%  sudah diperbaiki.
%=====================================================================
\section{QUANTUM}

\begin{frame}
    \frametitle{Prinsip Dasar}
    \begin{columns}[T]
        \begin{column}{0.5\textwidth}
            \begin{itemize}\tiny\justifying
                \item Quantum \\
                    Quantum adalah teori fisika yang menggambarkan perilaku
                    partikel sangat kecil (atom dan partikel subatom),
                    didasarkan pada sifat dualitas: partikel dapat berperilaku
                    sebagai gelombang dan partikel sekaligus. Konsep dasarnya
                    dijelaskan oleh vektor keadaan $\ket{\psi}$
                    \cite{barnett2009quantum}.
                \item Operator Unitary \\
                    Operator unitary digunakan dalam evolusi kuantum dan
                    pemrosesan informasi kuantum: ia menentukan apa yang
                    mungkin dan tidak mungkin dilakukan dalam pemrosesan
                    informasi kuantum,
                    \begin{equation}
                        \hat{U}\hat{U}^\dagger=\hat{I}=\hat{U}^\dagger\hat{U}.
                    \end{equation}
                \item Matriks Densitas \\
                    Representasi matematis keadaan kuantum sistem, mencakup
                    keadaan campuran (\textit{mixed state}) ketika terdapat
                    ketidakpastian pada keadaan sistem; vektor keadaan cukup
                    hanya untuk keadaan murni. Dengan matriks densitas kita
                    dapat menghitung nilai ekspektasi observabel kuantum.
            \end{itemize}
        \end{column}

        \begin{column}{0.5\textwidth}
            \vspace{0.4cm}
            \begin{itemize}\tiny
                \item[]\justifying Definisi umum matriks densitas untuk
                    keadaan campuran:
                    \begin{equation}
                        \rho=\sum_i p_i\,\ketbra{\psi_i}{\psi_i}
                        \label{eq:matriks-densitas}
                    \end{equation}
                    dengan $p_i$ peluang sistem berada pada keadaan murni
                    $\ket{\psi_i}$, di mana $\sum_i p_i=1$ dan
                    $\Tr(\rho)=1$.
                \item[]\vspace{0.2cm}
                \item[]\justifying {\footnotesize Tips notasi: proyektor
                    $\ketbra{\psi_i}{\psi_i}$ cukup ditulis dengan makro
                    \code{\ketbra{a}{b}} dari \code{ppt.sty}.}
            \end{itemize}
        \end{column}
    \end{columns}
\end{frame}

\begin{frame}
    \frametitle{QUANTUM ENTANGLEMENT \cite{muller2023}}
    \begin{columns}[T]
        \begin{column}{0.5\textwidth}
            \begin{itemize}\tiny\justifying
                \item EPR \\
                    Pada tahun 1935, Albert Einstein, Boris Podolsky, dan Nathan
                    Rosen menerbitkan makalah yang menggambarkan eksperimen
                    pemikiran (Gedankenexperiment) untuk mengilustrasikan
                    ketidakmasukakalan \textit{quantum entanglement} yang
                    menantang hukum dasar alam semesta \cite{einstein1935can}.
                \item Quantum Entanglement \\
                    Secara sederhana, \textit{quantum entanglement} merujuk pada
                    ketergantungan aspek-aspek sepasang partikel terjerat:
                    aspek satu partikel bergantung pada aspek partikel lainnya,
                    berapa pun jarak di antara keduanya.
                \item Bentuk Partikel \\
                    Partikel terjerat dapat berupa elektron atau foton;
                    aspeknya dapat berupa keadaan partikel, misalnya arah
                    spin (berputar ke satu arah atau arah lain).
                \item Pengukuran \\
                    Apabila kedua partikel terjerat, pengukuran pada satu
                    partikel langsung menentukan informasi pasangannya, bahkan
                    ketika keduanya terpisah jutaan tahun cahaya.
            \end{itemize}
        \end{column}

        \begin{column}{0.5\textwidth}
            \begin{itemize}\tiny\justifying
                \item Superposisi \\
                    Superposisi kuantum adalah gagasan bahwa partikel berada
                    dalam beberapa keadaan sekaligus; saat diukur, partikel
                    memilih salah satu keadaan dalam superposisi tersebut.
                \item Cara Partikel Terjerat \\
                    Prinsip membuat partikel terjerat: memecah sistem menjadi
                    dua dengan jumlah besaran diketahui. Contoh: memecah
                    partikel dengan spin nol menjadi dua partikel berspin
                    berlawanan sehingga jumlahnya tetap nol.
                \item Runtuhnya Probabilitas \\
                    Dalam mekanika kuantum, spin tiap partikel adalah
                    "sebagian naik dan sebagian turun" sampai diukur; saat
                    pengukuran terjadi, keadaan kuantum spin "runtuh" ke naik
                    atau turun seketika, meruntuhkan partikel lainnya ke spin
                    berlawanan.
                \item Variabel Tersembunyi \\
                    Einstein berteori adanya properti tak diketahui --- dijuluki
                    variabel tersembunyi --- yang menentukan keadaan partikel
                    sebelum pengukuran.
                \item Persamaan Bell \\
                    Bell menghasilkan \textit{Bell's inequality}: selalu benar
                    untuk teori variabel tersembunyi, tetapi tidak selalu benar
                    untuk mekanika kuantum; uji eksperimen mendukung mekanika
                    kuantum \cite{aspect1999bell}.
            \end{itemize}
        \end{column}
    \end{columns}
\end{frame}
